\subsection{Security-model separation and a concrete quantum obstruction}\label{sec:security14:models}
The soundness and privacy results above are
\emph{classical interactive/CRS} results.
Random-oracle compilation and quantum adversaries
change both the proof model and the hardness
assumptions. In particular, the present
Pedersen commitment instantiation is
\emph{not post-quantum binding}.

\begin{definition}[Four distinct claims]\label{def:security14:models}
We distinguish:
\begin{enumerate}[label=\textup{(\roman*)}]
\item \emph{Classical interactive soundness:}
the verifier samples every required
outer challenge after the prescribed
commitment/proof prefix, and the
relation engine satisfies its
knowledge-extraction contract
against classical PPT provers.
\item \emph{Classical online ZK in a simulatable CRS:}
the commitment setup may be generated
by a simulator retaining $\alpha$;
malicious verifiers and query indices
are arbitrary classical causal
algorithms, and relation proofs
satisfy adaptive multi-proof WI.
\item \emph{Fiat--Shamir ROM security:}
the interactive challenges are
derived from a domain-separated
transcript hash, and a separate
classical random-oracle proof
justifies every extraction and
simulation step in that
noninteractive experiment.
\item \emph{QROM/post-quantum security:}
the adversary may perform quantum
computation and quantum oracle
queries. All commitment,
proof-engine, hash and
folding-oracle reductions
must be valid against that
adversary class.
\end{enumerate}
Only (i)--(ii) are addressed by
the proved conditional
theorems of this chapter.
\end{definition}

\begin{lemma}[A known Pedersen setup logarithm gives two openings]\label{lem:security14:known-alpha}
If an adversary knows
$\alpha=\log_g h$, then for
every bit commitment
$C=g^b h^r$ it can efficiently
produce an opening to any
desired bit $b'\in\{0,1\}$:
\[
 \boxed{
 r'=r+(b-b')\alpha^{-1}\pmod P.}
\]
If $b'\neq b$ these are two
different valid bit openings.
\end{lemma}
\begin{proof}
Compute
\[
 g^{b'}h^{r'}
 =g^{b'}g^{\alpha r+b-b'}
 =g^b h^r=C.
\]
All operations are in
the public prime-order
scalar field; $\alpha$
is nonzero, so its inverse
exists. If $b'\ne b$,
the two openings have
different bit values
although the public
commitment is unchanged.
\end{proof}

\begin{theorem}[Quantum break of classical Pedersen-based binding]\label{thm:security14:quantum-break}
Assume the commitment group is a
standard finite cyclic group
whose discrete logarithm can
be computed by a polynomial-time
quantum algorithm
\cite{Shor1997}.
Then the field-symbol
commitments of Chapter~13
are \emph{not} quantum
computationally binding.
Indeed, a quantum adversary
can run the following
false-claim attack on the
full-relation protocol:
register a valid preprocessing
proof for $f_{\rm pre}\equiv0$,
wait for the evaluation point
$z$, and subsequently produce
valid proofs for an equivocated
source $f^{\rm post}\equiv1$
and public claim $v=1$.
All public commitment vectors
and the initial Merkle root
remain unchanged. Whenever the
underlying relation engine is
perfectly complete and allows
honest proofs of its true
statements, this attack
passes the algebraic and
authentication checks while
$1\neq\widetilde f_{\rm pre}(z)=0$.
\end{theorem}
\begin{proof}
First sample the honest
preprocessing commitment
arrays and all required
blinding scalars for the
zero source and zero mask.
Construct $Q_0=0$, its
full authenticated
commitment tree, and
an honest proof of the
true pre-evaluation
relation. Submit the
result before learning $z$.
The preprocessing source
is therefore the zero
table, so
$\widetilde f_{\rm pre}(z)=0$
for every $z$.

After learning $z$,
use the assumed quantum
discrete-log algorithm to
obtain
$\alpha=\log_g h$
from the already published
key, unless it was computed
earlier. For every previously
registered Pedersen bit
commitment, use
Lemma~\ref{lem:security14:known-alpha}
to compute an opening to
the bits of the desired
constant-one source,
zero mask and corresponding
masked $Q_0$. The public
commitment vectors, their
hashes, the first Merkle
root and its paths
do not change.
The pre-evaluation proof
remains a valid proof of
its original true
statement. All subsequent
statements are also true:
take the virtual
constant-one table,
generate its honest
Sumcheck masks and
fold chain, equivocate
the original level-zero
leaf commitments to
that chain, and build
the later levels
honestly.
The identities and
completeness of
Theorem~\ref{thm:full:completeness}
make every post-evaluation
NP relation and
authenticated query valid.
The adversary possesses
every required witness
and blinding and can
therefore produce honest
proofs without violating
their proof-system
soundness.

The verifier accepts
by perfect completeness,
yet the claimed value
is $v=1$ while the
preprocessing source
had evaluation $0$.
Finally, Shor's
quantum discrete-logarithm
algorithm~\cite{Shor1997}
runs in quantum
polynomial time on the
stated cyclic-group
families. Thus neither
the binding property
nor the protocol's
evaluation binding
survives this quantum
adversary.
\end{proof}

\begin{corollary}[No post-quantum theorem for the present Pedersen instantiation]\label{cor:security14:no-pq}
A claimed QROM or post-quantum
knowledge-soundness theorem
for the \emph{unchanged}
Pedersen-based full protocol
would be false under
the quantum group
model of
Theorem~\ref{thm:security14:quantum-break}.
\end{corollary}
\begin{proof}
The theorem explicitly
constructs a
polynomial-time quantum
false-claim adversary
accepted by the honest
verifier. Therefore the
required quantum
evaluation-binding game
does not have negligible
adversarial success.
A proof model cannot
restore a binding
property already
broken by the group's
quantum discrete-log
algorithm.
\end{proof}

\begin{proposition}[Original KBFold QROM soundness does not automatically transfer]\label{prop:security14:qrom-nontransfer}
The original KBFold
knowledge-soundness/QROM
theorems cannot be applied
as black-box end-to-end
security proofs of
the Chapter~13
Pedersen-hiding protocol.
\end{proposition}
\begin{proof}
The original scheme
authenticates field
codeword values with
a transparent Merkle/FRI
oracle and proves
proximity for the
specified unmasked code
and challenge schedule.
Chapter~13 changes
the leaves to Pedersen
bit-commitment vectors,
hides every field
opening behind
independent NP proofs,
changes the source
relation and degree
budgets to $D_j$,
and introduces
a distinct simulator
and extraction interface.
The assumptions and
transcript distribution
of the old theorem
therefore do not match
those of the new
protocol. Moreover
Theorem~\ref{thm:security14:quantum-break}
exhibits a quantum
binding failure in
the new leaf scheme,
independently of the
original hash/FRI
proximity argument.
Consequently no
automatic theorem
transfer is valid.
\end{proof}

\begin{definition}[Required post-quantum replacement contract]\label{def:security14:pq-replacement}
A future QROM secure
version of the \emph{same
Chapter~13 protocol
structure} must replace
the Pedersen binding layer
by a quantum-resistant
commitment with a
suitable equivocal
simulation setup or
another proved
online privacy method;
instantiate a
quantum-secure
relation argument
with compatible
extraction and adaptive
simulation;
and re-prove
authentication,
proximity and
noninteractive
challenge-sampling
bounds for the complete
altered transcript.
These are \emph{Chapter~15
certification/open-obligation
entries}, not a new
Chapter~16.
\end{definition}

\begin{remark}[Fiat--Shamir is not a missing algebraic lemma]
Hashing an interactive
verifier message into a
Fiat--Shamir challenge
changes the prover's
ability to correlate
its commitments and
challenge oracle
queries. Without an
appropriate ROM theorem
for the exact engine
and transcript order,
the classical interactive
proofs of
Theorems~\ref{thm:security14:full-binding}
and~\ref{thm:security14:online-zk}
do not imply
noninteractive security.
No such reduction is
claimed in this chapter.
\end{remark}
